\SourcePage{231}{236}
\section{Radiation of diatomic molecules. Vibrational and rotational spectra}

The selection rules and formulae for the transition probabilities listed in the preceding paragraph retain their validity also for transitions in which the electronic state of the molecule does not change\,\footnote{Transitions with change of vibrational (and with it rotational) state form, as one says, the \textit{vibrational spectrum} of the molecule; it lies in the near infrared region (wavelengths less than 20 $\mu$m). Transitions with change of the rotational state only form the \textit{rotational spectrum}, lying in the far infrared region (wavelengths greater than 20 $\mu$m).}. We shall dwell here only on some specific features of these transitions.

First of all, by the rule~(53.4) (dipole) transitions without change of the electronic state are altogether forbidden in molecules of identical atoms, since in such a transition the parity of the electronic term would remain unchanged. As follows from what was said in~\S\,53, this prohibition would be violated only when the interaction of nuclear spins with electrons is taken into account, while for molecules from different isotopes of one and the same element it is already violated by the influence of the rotation on the electronic state.

\SourcePage{232}{237}
The computation of matrix elements of the dipole moment reduces (according to formulae of~\S\,87 (see~III)) to their computation in the system of coordinates rotating with the molecule. The wave function of the molecule in this system is a product of the electronic wave function at a given distance $r$ between the nuclei and the wave function of the vibrational motion of the nuclei in the effective field $U(r)$ of the electrons and the nuclei. With complete neglect of the influence of the motion of the nuclei on the electronic state, the initial and final electronic wave functions in the transitions under consideration are the same. Integration over the electronic coordinates gives therefore in the matrix element simply the average dipole moment of the molecule $\bar{d}$ (directed, obviously, along the axis of the molecule)\,\footnote{In a molecule of identical atoms $\bar{d} = 0$, which is obvious from symmetry considerations.}. By virtue of the smallness of the vibrations the function $\bar{d}(r)$ can be expanded in powers of the vibrational coordinate $q = r - r_0$. For transitions involving a change of the vibrational state, the zeroth-order term of the expansion drops out of the matrix element by virtue of the orthogonality of the wave functions of the vibrational motion in one and the same field $U(q)$, so that there remains the term proportional to $q$. If the vibrations are treated as harmonic, then according to the well-known properties of the linear oscillator (see~III, \S\,23) the matrix elements are non-zero only for transitions between adjacent vibrational states; in other words, for the vibrational quantum number $v$ the selection rule
\begin{equation}
v' - v = \pm 1.
\tag{54.1}
\end{equation}
is valid. This rule is violated, however, when the anharmonicity of the vibrations is taken into account, as well as by the following terms of the expansion of the function $\bar{d}(q)$.

For a purely rotational transition (without change of the vibrational and electronic states) the matrix element of the projection of the dipole moment onto the moving axis $\zeta$ can be set simply equal to the average dipole moment of the molecule $\bar{d} = \bar{d}(0)$\,\footnote{In a molecule of identical atoms $\bar{d} = 0$, which is obvious from symmetry considerations.}. For the probability of the transition $J \to J \to 1$ we obtain as a result the formula
\begin{equation}
w(nJ \to n,\,J-1) = \frac{4\omega^3}{3\hbar c^3}\,\bar{d}^2\,\frac{J^2 - \Omega^2}{J(2J+1)},
\tag{54.2}
\end{equation}
which makes it possible to compute not only the relative (as in~(53.12)) but also the absolute values of the probabilities. (Formula~(54.2) is written for case $a$; in case $b$ one has to write $K$, $\Lambda$ instead of $J$, $\Omega$.)

The frequencies of purely rotational transitions are given by the differences of the rotational energies $BJ(J + 1)$ and equal
\begin{equation}
\hbar\omega_{J,J-1} = 2BJ.
\tag{54.3}
\end{equation}
The successive lines are equally spaced (at intervals $2B$) from each other.
